\documentclass[main.tex]{subfiles}


\begin{document}

%from pagenumber to up
% \phantomsection 
% \hypertarget{toc}{}

 

 

% номер секции вручную
\setcounter{section}{31
}
\addtocounter{section}{-1} 

\section{Работа. Простые механизмы}

\textbf{Работа} ($A$ [Дж])~--- это характеристика <<разгоняющего>> действия силы:
\begin{equation}
    A=Fr\cos\alpha,
\end{equation}
где~$r$~--- перемещение точки приложения силы (cтрого говоря, $\Delta r$), $\alpha$~--- угол между векторами силы и перемещения. (Если направление силы в процессе движения меняется, но угол~$\alpha$ на малых перемещениях остается постоянным, то работа равна: $A=FS\cos\alpha$, где $S$~--- путь точки приложения силы.)

\textbf{Мощность} ($N$ [Вт])~--- это быстрота совершения работы:
\begin{equation}
    N=\frac{A}{t}.
\end{equation}
%
% Например, при свободном падении тяжелого шара из состояния покоя у поверхности планеты сила тяжести, приложенная к нему, совершает положительную работу, что ведет к увеличению скорости шара.

\textbf{Простым механизмом} называют механическое устройство для преобразования силы. На рис.\,\ref{fig:p56} представлены примеры таких механизмов.

\begin{figure}[H]
    \centering

    \begin{tikzpicture}[font=\small,            line cap=round,                             >={Stealth[inset=0pt,round,length=3mm, angle'=20]},vec/.style={>={Stealth[inset=0pt,round,length=3.5mm, angle'=20]}}]


% \filldraw[lightgray,semitransparent] (0,0) -- ({6-cos(45)*.5},0) --++ (45:.5) --++ (-45:.5) -- ({9-cos(45)*.5},0) --++ (45:.2) coordinate (P) arc [start angle=-135, end angle=-90, radius=.3]
% \filldraw[lightgray,semitransparent] (0,0) -- ({6-cos(45)*.5},0) --++ (45:.5) --++ (-45:.5) -- ({10-cos(45)*.5},0) --++ (45:.5) coordinate (P) --++ (0,{-cos(45)*.5-.3}) -- (0,-.3) -- cycle;
% \draw (0,0) -- ({6-cos(45)*.5},0) --++ (45:.5) --++ (-45:.5) -- ({10-cos(45)*.5},0) --++ (45:.5) --++ (0,{-cos(45)*.5-.3});
\filldraw[lightgray,semitransparent] (0,0) -- ({6-cos(45)*.5},0) --++ (45:.5) --++ (-45:.5) -- ({10-cos(60)*3},0) --++ (60:3) coordinate (P) --++ (0,{-.3/cos(60)}) -- ({10-cos(60)*3+.3*tan(60/2)},-.3) -- (0,-.3) -- cycle;
\draw (0,0) -- ({6-cos(45)*.5},0) --++ (45:.5) --++ (-45:.5) -- ({10-cos(60)*3},0) --++ (60:3) coordinate (P) --++ (0,{-.3/cos(60)});

%%%%%%%%%%%%%%%%%%%%%%%%%%%%%%
%%%%plane pic

%calc center of cube
\path[name path=g] (7,{sin(45)*.5+.1+.3}) --++ (-7,0);
\path[name path=dc] (1,0) --++ (75:1);
\path[name intersections={of=g and dc, by=int}] (int) coordinate (ccube) {};
\filldraw[lightgray,draw=black,rotate around={30:(1,0)}] let \p1 = ($ (int) - (1,0) $) in (1,0) rectangle++ ({sqrt(2)*sqrt(\x1^2+\y1^2)},{sqrt(2)*sqrt(\x1^2+\y1^2)});

%plane
\fill[brown,semitransparent] (1,0) --++ (30:2.5) --++ (0,-1.25) -- cycle;
\draw (1,0) --++ (30:2.5) node [black,above] {П} --++ (0,-1.25) -- cycle;
\filldraw[lightgray,draw=black,rotate around={30:(1,0)}] let \p1 = ($ (int) - (1,0) $) in (1,0) rectangle++ ({sqrt(2)*sqrt(\x1^2+\y1^2)},{sqrt(2)*sqrt(\x1^2+\y1^2)}); % calculated above (solid)
% \filldraw[lightgray,draw=black,rotate around={30:(1,0)}] (1,0) rectangle++ (1,.5);

\draw[->,red,thick,vec] (ccube) coordinate (F)--++ (-90:2.2) node[black,right,inner ysep=0] {$\vec F_\text{т}$};
\pgfgetlastxy{\ix}{\iy}
\draw[->,Green,thick,vec] (ccube) coordinate (F)--++ (30:1.1) node[black,above] {$\vec F$};
\node[] at (F) [circle,fill=Green,inner sep=0pt,minimum size=1mm,label=] {};

%lever
\fill[brown,semitransparent] (4,{sin(45)*.5}) rectangle (7,{sin(45)*.5+.1});
\draw[] (4,{sin(45)*.5}) rectangle (7,{sin(45)*.5+.1});
\node [black,above] at (5.5,{sin(45)*.5+.1}) {Р};

\shade[ball color=lightgray] (7,{sin(45)*.5+.1+.3}) coordinate (C) circle (.3);
\path (C);
\pgfgetlastxy{\cx}{\cy}

\draw[->,red,thick,vec] (C) coordinate (F) -- (7,\iy) node[black,right,inner ysep=0] {$\vec F_\text{т}$};
\node[] at (F) [circle,fill=red,inner sep=0pt,minimum size=1mm,label=] {};

\draw[->,Green,thick,vec] let \p1 = ($ (C) - (7,\iy) $) in (4,{sin(45)*.5+.05}) coordinate (F)--++ (0,{-\y1/2}) node[black,right] {$\vec F$};
\node[] at (F) [circle,fill=Green,inner sep=0pt,minimum size=1mm,label=] {};

%%%%%%%%%%%%%%%%%%%%%%%%%%%%%%%%%%%%%%
%%%% pulley

\fill[brown,semitransparent] (P) --++ (-90:.3) arc [start angle=-90, end angle=240, radius=.3] -- cycle;
\draw[] (P) ++ (-90:.3) arc [start angle=-90, end angle=240, radius=.3];
\node[] at (P) [circle,fill=black,inner sep=0pt,minimum size=1mm,label=] {};
\node[above] at ([yshift=.3cm]P) {Б$_\text{н}$};

%thread
\path (P);
\pgfgetlastxy{\px}{\py}
\draw[] (P) ++ (150:.3) ++ ({150+90}:.5) coordinate (F) --++ ({150+90-180}:.5) arc [start angle=-210, end angle=-360, radius=.3] -- ({\px+.3cm},{\cy+.3cm});

%solid
\shade[ball color=lightgray] ({\px+.3cm},{\cy}) coordinate (C) circle (.3);


%F
\draw[->,Green,thick,vec] let \p1 = ($ (C) - (7,\iy) $) in (F) --++ ({150+90}:\y1) node[black,left] {$\vec F$};
\node[] at (F) [circle,fill=Green,inner sep=0pt,minimum size=1mm,label=] {};

\draw[->,red,thick,vec] (C) -- ({\px+.3cm},{\iy}) node[black,right,inner ysep=0] {$\vec F_\text{т}$};
\node[] at (C) [circle,fill=red,inner sep=0pt,minimum size=1mm,label=] {};


%%%%%%%%%%%%%%%%%%%%%%%%%%%%%%%%%%%%%%
%%%% movable pulley 
\def\dpul{13.6} % 3rd pulley coord for equal distance
\filldraw[lightgray,semitransparent] (\dpul,4) rectangle (15,4.3);
\draw (\dpul,4.3) --++ (0,-.3) -- (15,4);

%circle pulley
\fill[brown,semitransparent] (\dpul,\py) coordinate (P) circle (.3);
\draw[] (P)  circle (.3);
\node[] at (P) [circle,fill=black,inner sep=0pt,minimum size=1mm,label=] {};
\node[right] at ([xshift=.3cm]P) {Б$_\text{п}$};


%thread
\path (P);
\pgfgetlastxy{\px}{\py}
\draw[] (P) ++ (180:.3) ++ ({90}:.5) coordinate (F) --++ ({-90}:.5) arc [start angle=180, end angle=360, radius=.3] -- ({\dpul+.3},4);
\draw (P) -- ({\px},{\cy+.3});

%solid
\shade[ball color=lightgray] ({\px},{\cy}) coordinate (C) circle (.3);

%F
\draw[->,Green,thick,vec] let \p1 = ($ (C) - (7,\iy) $) in (F) --++ ({90}:{\y1/2}) node[black,left,inner ysep=0] {$\vec F$};
\node[] at (F) [circle,fill=Green,inner sep=0pt,minimum size=1mm,label=] {};

\draw[->,red,thick,vec] (C) -- ({\px},{\iy}) node[black,right,inner ysep=0] {$\vec F_\text{т}$};
\node[] at (C) [circle,fill=red,inner sep=0pt,minimum size=1mm,label=] {};



    \end{tikzpicture}
\caption{Простые механизмы}
\label{fig:p56}
\end{figure}

Простые механизмы (обозначены светло"=коричневым цветом) используются для совершения \emph{полезного процесса} (в данном случае подъема массивных куба и шаров). Зеленый вектор~$\vec F$~--- сила, приложенная к механизму\footnote{В случае \emph{наклонной плоскости} силу прикладывают к телу, установленному на нее.} (\emph{затраченная сила}); красный вектор~$\vec F_\text{т}$~--- сила тяжести, действующая на груз. 
% (остальные силы, действующие на тела, не указаны).
% Вот основные виды простых механизмов.  
\begin{enumerate}
    \item Гладкая закрепленная \textbf{наклонная плоскость}~П позволяет получить выигрыш в силе: $F<F_\text{т}$. Меняется также направление необходимого усилия.
    \item Невесомый \textbf{рычаг}~Р~--- это твердое тело, способное вращаться вокруг неподвижной точки в одной плоскости. Он также позволяет получить выигрыш в силе: $F<F_\text{т}$. Направление необходимого усилия меняется.
    \item Гладкий \textbf{неподвижный блок}~Б$_\text{н}$~--- колесо с желобом~--- с пропущенной по желобу нитью выигрыша в силе не дает ($F=F_\text{т}$), но этот блок позволяет изменить направление прикладываемого усилия.
    \item Невесомый \textbf{подвижный блок}~Б$_\text{п}$~--- колесо с желобом, способное перемещаться в своей плоскости,~--- с помощью двух нитей в наиболее выгодном положении (рис.\,\ref{fig:p56}, справа) дает выигрыш в силе ($F<F_\text{т}$).
    
\end{enumerate}















 
\end{document}